\chapter{Introduction}
\label{ch:introduction}

\section{Motivation}

The increasing demand for autonomous logistics and the growth of large-scale goods transportation have positioned articulated tractor-trailer vehicles as a critical target for autonomous control research. The transition from rigid-body autonomous navigation to the control of articulated tractor-trailer systems represents a significant escalation in mechanical complexity, perception requirements, and control instability, particularly during reverse maneuvering.

Unlike single-body vehicles, tractor-trailer systems introduce non-holonomic constraints at the hitch point that fundamentally alter the stability characteristics of the vehicle. In reverse motion, the articulation angle $\gamma$ between the tractor and trailer exhibits open-loop instability: absent precise closed-loop control, any perturbation grows unbounded, culminating in the well-known jackknife condition from which no steering correction can recover the trailer's alignment. Classical approaches such as Nonlinear Model Predictive Control (NMPC) can manage this instability by optimising over a prediction horizon, but their computational demands and reliance on accurate system models limit their applicability in unstructured or sensor-rich environments.

\section{Problem Definition}

The core challenge addressed in this thesis is twofold. First, the kinematic and dynamic complexity of tractor-trailer systems makes it difficult to design reward functions and control policies that simultaneously optimise path tracking performance and hitch stability. Second, real-world autonomous deployment requires perception of lane boundaries and obstacles directly from high-dimensional sensor data, rather than from pre-computed symbolic state representations.

Existing learning-based path-following studies on rigid-body vehicles have shown that continuous-control reinforcement learning can achieve performance competitive with classical baselines under nominal conditions, but robustness degrades when the vehicle model, operating conditions, or path geometry become more demanding. Those limitations are amplified for articulated systems, where reverse instability and hitch-angle growth impose a much narrower margin for control error.

\section{Thesis Contributions}

This thesis addresses these challenges through the following contributions:

\begin{itemize}
    \item developing a kinematic and dynamic model of a tractor-trailer system, including hitch instability analysis for reverse manoeuvring;
    \item designing a vision-based RL observation space using CNN encoders and autoencoders operating on bird's-eye-view (BEV) images rendered from a custom Pygame simulation environment;
    \item engineering a reward function that balances cross-track error, heading error, hitch articulation stability, and collision avoidance;
    \item implementing and benchmarking Twin Delayed Deep Deterministic Policy Gradient (TD3) agents in lane-following and obstacle-avoidance tasks within a custom Gymnasium-based simulation environment;
    \item comparing the learned policies against tuned classical control baselines, including pure-pursuit, PID, and MPC, on matched articulated path-tracking tasks.
\end{itemize}

The research aligns with the objectives of the Mechatronics Vehicle Systems (MVS) Lab and the WATonoBus project at the University of Waterloo.

\section{Institutional Context}

This work is submitted in partial fulfilment of the requirements for a Master of Applied Science (MASc) degree in the Department of Mechanical and Mechatronics Engineering at the University of Waterloo. The degree requires six terms of full-time registration, completion of 3--5 graduate-level courses, appointment of two additional faculty readers, and a mandatory 15-day public display period. The target completion date is early July, with a seminar scheduled for late July.

\begin{table}[H]
\centering
\caption{MASc thesis requirements summary}
\label{tab:masc_requirements}
\begin{tabular}{ll}
\toprule
Element & Requirement \\
\midrule
Research Scope   & Substantial problem extending beyond existing laboratory results \\
Literature Review & Comprehensive assessment of classical control and modern AI approaches \\
Validation       & Simulation-based empirical evidence across multiple environment configurations \\
Length/Format    & Typically 80--120 pages; GSPA formatting and UWSpace submission \\
Committee        & Supervisor plus two faculty readers; formal display period \\
\bottomrule
\end{tabular}
\end{table}

\section{Thesis Organisation}

The remainder of this thesis is structured as follows. Chapter~\ref{ch:literature_review} reviews classical and learning-based control methods for articulated vehicles, alongside relevant perception and safety literature. Chapter~\ref{ch:modeling} derives the state-space equations for the tractor-trailer system and analyses the jackknife instability. Chapter~\ref{ch:perception} presents the CNN and autoencoder-based perception stack. Chapter~\ref{ch:rl_framework} defines the MDP formulation and details the TD3 learning framework used in this work. Chapter~\ref{ch:simulation} describes the custom Pygame/Gymnasium simulation environment and experimental setup. Chapter~\ref{ch:results} presents and discusses experimental results. Chapter~\ref{ch:conclusion} summarises the findings and outlines future work.
